% HD-N18. GPT-6 Astra Pro, /root/imprimitive; inherited effort.
% 21 November 2026. Hodge-Clifford positivity obstruction for ample secants.
\subsection{HD-N18: ample sixfold secants cannot become whole Weil characters}
\label{sec:HD-N18}

\paragraph{The precise class of actual objects.}
Consider actual perfect complexes on complex abelian sixfolds whose
whole character is
\begin{equation}
 e^{B_0}\bigl(\alpha e^{ic\theta_0}
                    +\beta e^{-ic\theta_0}\bigr),
 \qquad c>0,\quad \alpha\beta\ne0,                    \tag{N18.1}
\end{equation}
where $\theta_0$ is an ample real $(1,1)$ class, $B_0$ is a
real $(1,1)$ class, and the displayed character is rational.
Their full actual Fourier--Mukai, isogeny, arbitrary line-twist,
translation, shift and derived-dual orbits contain no whole character
\begin{equation}
 e^B\bigl(P(\theta)+\lambda t\bigr),\qquad
 B\in\operatorname{NS}(A)_{\mathbf Q},\quad
 \deg P\le6,\quad 0\ne\lambda t\in W_K(A),            \tag{N18.2}
\end{equation}
on a polarized Weil sixfold of signature $(3,3)$, for any imaginary
quadratic field $K$. Nonzero scalar multiples of transported
characters are included.

This is a whole-character obstruction for the specified actual
objects and functors. In particular it applies to any actual
rank-two locally free bundle with $c_1=0$ and
$c_2=d\theta_0^2$, $d>0$, since its exact rank-two Chern-root
identity gives $\operatorname{ch}(E)=e^{i\sqrt d\theta_0}
+e^{-i\sqrt d\theta_0}$. No existence assertion for such a bundle
is used. The obstruction closes its proposed transport as a Weil
supplier even before an existence or semiregularity construction.
It concerns neither a different character nor a different
correspondence outside the listed actual functorial orbit.

\paragraph{1. The exact Hodge Clifford space and its functor maps.}
For an abelian variety $A$ let $V=H^1(A,\mathbf R)$, and let
$I$ be its real Hodge complex structure, acting by $i$ on
$H^{1,0}$ and by $-i$ on $H^{0,1}$. Put
\begin{equation}
 \mathcal W=V\oplus V^*,\quad
 Q(v+\xi,w+\eta)=\xi(w)+\eta(v),\quad
 \mathcal I=\operatorname{diag}(I,-I^*),\quad
 c(v+\xi)=v\wedge+\iota_\xi.                         \tag{N18.3}
\end{equation}
Here $I^*\xi=\xi\circ I$. The $-i$ eigenspace of
$\mathcal I$ is the actual Hochschild-generator space
\[
 U=H^{0,1}(A)\oplus H^{1,0}(A)^*
                       =H^1(\mathcal O_A)\oplus H^0(T_A).
\]
It is $Q$-isotropic, and $HH^2(A)=\bigwedge^2U$ acts by
the corresponding products of these Clifford operators.

For an actual Fourier--Mukai equivalence, the intertwiner of
HD-G19 and HD-N17 is a real isometry $f$ with
$Lc(w)L^{-1}=c(fw)$. It also intertwines $\mathcal I$.
Here is the sign identification: the Poincar\'e swap
$\iota_A(v,\xi)=(\xi,v)$ identifies $\mathcal W$ with
$H_1(A)\oplus H_1(\widehat A)$. The actual holomorphic tangent
space of $A\times\widehat A$ is
$H^{1,0}(A)^*\oplus H^{0,1}(A)$, which is the swap of $U$.
The graph isomorphism of the equivalence is holomorphic, so its
map carries this space to the analogous target space. Equivalently
$f\mathcal I_A=\mathcal I_{A'}f$; the complex structure of the
point variety corresponds to $-\mathcal I$, and the same minus
sign occurs on both sides. This is precisely the Hodge refinement
of the original GLO Proposition4.3.7/Orlov Theorem2.10 map used
in HD-G20 and HD-R05, not an assumption on an arbitrary
orthogonal transformation.

An isogeny uses $f=(T,(T^{-1})^*)$, $T=p^*$ on $H^1$,
and is Hodge linear. Its pushforward is a nonzero scalar times
the inverse cohomological map. Multiplication by $e^D$ uses
the literal shear
\begin{equation}
       f_D(v,\xi)=(v-\iota_\xi D,\xi).               \tag{N18.4}
\end{equation}
For $D$ of type $(1,1)$, its contraction map obeys
$ID^\sharp+D^\sharp I^*=0$, so (N18.4) commutes with
$\mathcal I$. It is a $Q$-isometry. These facts apply to
arbitrary mixed line classes and also to the real $(1,1)$
shears removing $B_0$ and $B$ in the character calculations.
Removing these displayed factors is an explicit linear map;
no geometric line bundle with a prescribed real class is asserted.

On even spinors, derived dual is
$\tau|_{\wedge^j V}=(-1)^{j(j-1)/2}$. Its action on
annihilators is the explicit anti-isometry
\begin{equation}
          f_{\mathrm{dual}}(v,\xi)=(v,-\xi).          \tag{N18.5}
\end{equation}
Indeed for an even spinor $u$,
$\tau c(w)u=c(f_{\mathrm{dual}}w)\tau u$;
the two parity calculations are that a wedge has the same sign
and a contraction has the opposite sign. On an odd spinor the
formula has an additional minus sign. Consequently on degree-two
Hochschild operators the conjugation is
$-c(f_{\mathrm{dual}}w_1)c(f_{\mathrm{dual}}w_2)$.
This preserves their ranks. The map (N18.5) commutes with
$\mathcal I$ and multiplies $Q$ by $-1$. It conjugates the
two bivector summands of the spin Lie algebra by a minus sign
and fixes its middle $\mathfrak{gl}(V)$ summand, exactly as
in HD-N17. Translations and shifts have the identity map on
$\mathcal W$, with respectively the identity or a scalar on
characters. Thus every permitted composition gives a real
isometry or anti-isometry intertwining the Hodge complex
structures, preserving Hochschild ranks, carrying pure-spinor
annihilators, and conjugating full spin Lie stabilizers.

\paragraph{2. Rank thirty forces a pure target Weil character.}
Remove $e^{B_0}$ by (N18.4). On the two pure exponentials in
(N18.1), an element $(\eta,v)\in
H^{0,1}\oplus H^{1,0*}$ acts through the one-forms
\[
              \pi_\pm(\eta,v)=\eta\pm ic\,\iota_v\theta_0.
\]
Since $\theta_0$ is nondegenerate and $c\ne0$, the actual
map $(\pi_+,\pi_-):U\to H^{0,1}\oplus H^{0,1}$ is an
isomorphism: its sum recovers $2\eta$ and its difference
recovers $2ic\,\iota_v\theta_0$. In the resulting two
six-dimensional graph summands of $U$, mixed degree-two
products kill both exponentials, while the two exterior
squares have images
\[
        \bigwedge^2H^{0,1}\,e^{ic\theta_0},
             \qquad \bigwedge^2H^{0,1}\,e^{-ic\theta_0}.
\]
Each has dimension fifteen. Their intersection is zero:
an equality $\eta e^{ic\theta_0}=\zeta e^{-ic\theta_0}$
gives $\eta=\zeta$ in degree two and
$2ic\theta_0\eta=0$ in degree four. The map
$\theta_0:H^{0,2}\to H^{1,3}$ is injective. In a diagonal
positive $(1,1)$ basis, the coefficient of an antiholomorphic
pair is recovered by wedging with any third, distinct paired
index. This proves injectivity directly. Nonzero $\alpha,\beta$
do not change either image. Therefore the exact source rank is
\begin{equation}
              \operatorname{rank}c^2_{\operatorname{ch}(E)}=30.
                                                        \tag{N18.6}
\end{equation}

For every target (N18.2), the full sixfold calculation in
HD-R06 gives
\[
 \operatorname{rank}c^2=30+9\operatorname{rank}M_P,
 \qquad (M_P)_{ij}=(i+j)!p_{i+j},\quad
                 0\le i\le4,\ 0\le j\le2,
\]
where $P(z)=\sum_{k=0}^6p_kz^k$.
The actual maps in paragraph1 preserve this rank. Thus every
entry of $M_P$ is zero and $P=0$. After the explicit target
shear, the whole target character is a nonzero rational Weil
class, a sum of its two nonzero conjugate determinant volumes.

\paragraph{3. Recovering the two pure axes from the full stabilizer.}
This step supplies the exact comparison of the two pairs;
no secant-decomposition uniqueness theorem is assumed.
On the source, the annihilators of
$\psi_\pm=e^{\pm ic\theta_0}$ are
\begin{equation}
 L_\pm^{\mathrm{src}}=
 \{\mp ic\,\theta_0^\sharp\xi+\xi:\xi\in V_{\mathbf C}^*\}.
                                                        \tag{N18.7}
\end{equation}
They are transverse maximal isotropic spaces of dimension twelve.
The actual relative Fock map of HD-N17, now in dimension twelve,
sends the vacuum and top exterior lines to $\psi_+$ and
$\psi_-$. The same triangular verification given there applies
with $2sH$ replaced by $ic\theta_0$.
For $a\cdot1+b\cdot\mathrm{vol}$ with $ab\ne0$, the
degree-two and degree-ten outputs of a spin Lie operator force
its wedge and contraction bivectors to be zero. Its degree-zero
and top-degree equations force its middle matrix to have trace
zero. Its full complex stabilizer is exactly
$\mathfrak{sl}_{12}$ acting on this Fock space.

The fixed spinors of this Lie algebra form precisely the
vacuum--top plane. Indeed in each intermediate exterior degree
a basis monomial contains some index and omits another;
the corresponding diagonal trace-zero matrix has a nonzero
weight on it. No such monomial is invariant under all diagonal
trace-zero matrices. Finally the only pure lines in the fixed
plane are its two axes. A combination with both coefficients
nonzero has no Clifford annihilator: the two halves of its
degree-one action have independent Fock degrees one and eleven,
and each is injective. Each axis, by contrast, has its displayed
twelve-dimensional annihilator.

For the target pure Weil class, let
$V_{\mathbf C}=V_+\oplus V_-$ be its field eigenspaces,
each of dimension six. Its determinant volumes have annihilators
\begin{equation}
 L_+^{\mathrm{tar}}=V_+\oplus V_-^*,\qquad
 L_-^{\mathrm{tar}}=V_-\oplus V_+^*.                  \tag{N18.8}
\end{equation}
Their relative Fock map is explicit: its operators wedge in
one eigenspace basis and contract out the other. It gives the
same vacuum--top model in dimension twelve, as checked in
HD-R06. Thus the target full stabilizer has the same exact
fixed plane and only those two pure axes.

The actual cohomological transport, including the two displayed
$(1,1)$ shears, conjugates the full stabilizers. It therefore
carries their fixed planes to one another. It carries pure
spinors to pure spinors by the annihilator maps of paragraph1.
It follows that its real isometry or anti-isometry $f$ carries
the unordered pair (N18.7) to the unordered pair (N18.8).
This is the required exact axis comparison.

\paragraph{4. The source symmetric form is definite.}
Write $C=\theta_0^\sharp:V^*\to V$ and define the real
orthogonal complex structure
\begin{equation}
 \mathcal J_{\mathrm{src}}=
       \begin{pmatrix}0&cC\\-c^{-1}C^{-1}&0\end{pmatrix}.
                                                        \tag{N18.9}
\end{equation}
It squares to $-1$, is skew-adjoint for $Q$, and its
$\pm i$ eigenspaces are exactly (N18.7). The $(1,1)$
identity $IC+CI^*=0$ proves that it commutes with
$\mathcal I$. Hence
\begin{equation}
          G_{\mathrm{src}}(x,y)
              =Q(x,\mathcal I\mathcal J_{\mathrm{src}}y)
                                                        \tag{N18.10}
\end{equation}
is a symmetric nondegenerate real form.

Its positivity, with the chosen signs, is explicit. Choose
real one-forms $a_j,b_j$ with $Ia_j=-b_j$, $Ib_j=a_j$ and
$\theta_0=\sum_{j=1}^6h_j a_j\wedge b_j$, $h_j>0$.
Let $a_j^*,b_j^*$ be the actual dual basis. Then
\[
 \begin{array}{ll}
 \mathcal J_{\mathrm{src}}a_j=(ch_j)^{-1}b_j^*,&
 \mathcal J_{\mathrm{src}}b_j=-(ch_j)^{-1}a_j^*,\\
 \mathcal J_{\mathrm{src}}a_j^*=ch_jb_j,&
 \mathcal J_{\mathrm{src}}b_j^*=-ch_ja_j.
 \end{array}
\]
Since $\mathcal I=\operatorname{diag}(I,-I^*)$, these
formulas give a diagonal matrix for (N18.10), with entries
$(ch_j)^{-1}$ on $a_j,b_j$ and $ch_j$ on $a_j^*,b_j^*$.
Every entry is positive. Replacing
$\mathcal J_{\mathrm{src}}$ by its negative makes the form
negative definite, so its unordered pair of eigenspaces
still determines definiteness up to sign.

\paragraph{5. The target form is neutral, giving the contradiction.}
Normalize the field action on target $V$ to a real complex
structure $K_V$, with eigenspaces $V_+,V_-$. The real
orthogonal complex structure determined by (N18.8) is
\begin{equation}
 \mathcal J_{\mathrm{tar}}=\operatorname{diag}(K_V,-K_V^*).
                                                        \tag{N18.11}
\end{equation}
It commutes with $\mathcal I$ because the actual field
endomorphisms are Hodge endomorphisms. Its associated real
symmetric form is
\[
       G_{\mathrm{tar}}(x,y)
                      =Q(x,\mathcal I\mathcal J_{\mathrm{tar}}y).
\]
It is nondegenerate, and its restriction to $V$ is zero:
both complex structures preserve $V$, and $Q|_{V\times V}=0$.
The subspace $V$ has real dimension twelve inside the
twenty-four-dimensional $\mathcal W$. Thus the real
signature of this symmetric form is $(12,12)$; each of
the positive and negative indices is at least twelve
because of that half-dimensional isotropic subspace.

The exact maps in paragraphs1 and3 give
\[
 f\mathcal J_{\mathrm{src}}f^{-1}
       =\varepsilon\mathcal J_{\mathrm{tar}},\qquad
 f\mathcal I_{\mathrm{src}}=\mathcal I_{\mathrm{tar}}f,\qquad
 Q_{\mathrm{tar}}(fx,fy)=\epsilon Q_{\mathrm{src}}(x,y),
 \quad\varepsilon,\epsilon\in\{1,-1\}.
\]
Therefore the actual symmetric
forms obey
\begin{equation}
       G_{\mathrm{tar}}(fx,fy)
             =\varepsilon\epsilon\,G_{\mathrm{src}}(x,y).
                                                        \tag{N18.12}
\end{equation}
The right side is definite; the left side has signature
$(12,12)$. An invertible real map cannot satisfy this
identity. This proves the whole-character exclusion.

\paragraph{Scope of the completed source test.}
All steps apply to actual functors, with their explicit
Hodge maps and their exact even-spinor duality signs.
The proof establishes a closed obstruction for the ample
two-imaginary-exponential sixfold source class (N18.1).
It neither assumes existence of an unknown rank-two bundle
nor labels its possible bundle constructions impossible.
The proposed supplier orbit is excluded by the proved
rank, axis and symmetric-form comparisons.
